\begin{Answer}{4.7}
\begin{code}
double areaSquare(double s) { return s*s; }
\end{code}
\end{Answer}
\begin{Answer}{4.10}
It does not work.  To see why,
notice that if we pass in $a$ and $b$, then $x=a$ and $y=b$ at the beginning.
After the first line, $x=b$ and $y=b$.
After the second line $x=b$ and $y=b$.
The problem is that the first line overwrites the value stored in $x$ ($a$), and we can't recover it.
\end{Answer}
\begin{Answer}{4.13}
Either $-1$ or $3$ are possible answers if we are uncertain whether it will return a positive
or negative answer.  But we know it is one of these.  It won't be -5, for instance.
\end{Answer}
\begin{Answer}{4.14}
\begin{solevalenum}
\item
This solution is both incorrect and a bit confusing.
The phrase `both sides' is confusing--both sides of what?
We don't have an equation in this problem.
But there is a more serious problem.
If you thought it was correct, go back and try to figure out why it is incorrect before you
continue reading this solution.
The main problem is that although this may return
a value in the correct range, it doesn't always return the correct value.  In fact, what if
$(a\pmod{b} + b-1)$ is odd?  Mathematically, this would result in a non-integer result which is
clearly incorrect.  In most programming languages it would at least truncate and
return an integer--but again, not always the correct one.
This person focused on the wrong thing--getting the number in a particular range.  Although that
is important, they needed to think more about how to get the {\em correct} number.
They should have plugged in a few more values to double-check their logic.
\item
Incorrect. Generally speaking, $a\bmod b \neq -a\bmod b$.
In other words, returning the absolute value when the result is negative is almost always incorrect.
\item
Incorrect.  If you think about if for a few minutes you should see that this is just
one way of implementing the idea from the previous solution.
\item
Incorrect.  If $(a\bmod b)$ is negative, performing another $\bmod$ will still leave it negative.
\end{solevalenum}
\end{Answer}
\begin{Answer}{4.15}
There are several possible answers, but the slickest is probably:
\verb#(b+(a mod b)) mod b#.  Try it with both positive and negative
numbers for $a$ and convince yourself that it is correct.
\end{Answer}
\begin{Answer}{4.16}
\begin{solevalenum}
\item
This does not work.  What happens when $x=3.508$, for instance?
\item
This is incorrect for two reasons.  First, $1/2=0$ in most programming languages, so this
will {\em always} round down.
Second, even if we replaced this with $.5$ or $1.0/2.0$, it would round {\em up} at $.5$.
\item
Nice try, but still no good.  What if $x=2.5$?  This will round {\em up} to $3$.
Worse, what if $x=2.0001$?  Again, it rounds {\em up} to $3$ which is really bad.
\item
This one is correct.  Plug in values like $2$, $2.1$, and $2.5$ to see that it rounds down to
$2$ and values like $2.51$, $2.7$, and $2.9$ to see that it rounds up to $3$.
\end{solevalenum}
\end{Answer}
\begin{Answer}{4.17}
(a) 0; (b) 1; (c) 1; (d) 1; (e) 1 (f) 1; (g) 2; (h) 9; (i) 0; (j) -1; (k) -1; (l) -2.
\end{Answer}
\begin{Answer}{4.18}
\begin{solevalenum}
\item
$0.5$ is not an integer, and the \verb+floor+ function is not allowed.
\item
The \verb+floor+ function is not allowed.  Even if it were, this solution
doesn't work. $1/2$ is evaluated to $0$ so it doesn't help.
\item
This one works, but $0.5$ is not allowed so it does not follow the directions.
\end{solevalenum}
\end{Answer}
\begin{Answer}{4.19}
Two reasonable solutions include \verb@(n+m/2)/m@ and \verb@(2n+m)/(2m)@.
\end{Answer}
\begin{Answer}{4.22}
Here is one possibility:
\begin{code}
    int max(int x, int y, int z) {
        int w = max(x,y);
        return max(w,z);
    }
\end{code}
We will use a proof by cases.
\begin{itemize}
\item
If $x$ is the maximum, then $w=$\verb+max+$(x,y) = x$.
so it returns \verb+max+$(w,z)=x$, which is correct.
\item
If $y$ is the maximum, the argument is essentially the same the previous case.
\item
If $z$ is the maximum, then $w$ is either $x$ or $y$, but in either case $w\leq z$,
so it returns \verb+max+$(w,z)=z$.
\end{itemize}
In each case the algorithm returns the correct answer.
\end{Answer}
\begin{Answer}{4.25}
Here is a possible answer.
\begin{code}
    void HelloGoodbye(int x) {
        if(x >= 4) {
            if(x <= 6) {
                print("Hello");
            } else {
                print("Goodbye");
            }
        } else {
            print("Goodbye");
        }
    }
\end{code}
\end{Answer}
\begin{Answer}{4.26}
It is possible.  If you thought it wasn't, go back and try to write the algorithm before
reading any further.

Here is one way to do it using an extra variable and an additional conditional statement.
\begin{code}
    void HelloGoodbye(int x) {
        boolean sayGoodbye = true;
        if(x >= 4) {
            if(x <= 6) {
                sayGoodbye = false;
            }
        }
        if(sayGoodbye) {
            print("Goodbye");
        } else {
            print("Hello");
        }
    }
\end{code}
It can also be done by using a return statement
(this version is not recommended!):
\begin{code}
    void HelloGoodbye(int x) {
        if(x >= 4) {
            if(x <= 6) {
                print("Hello");
                return;
            }
        }
        print("Goodbye");
    }
\end{code}
The second solution is simpler,
but this sort of code (with somewhat random \verb+return+ statements in the middle
of them) can be tricky to debug if it is changed later.
Did you come up with a better solution than these?
\end{Answer}
\begin{Answer}{4.27}
\verb+list.size()!=0+ and \verb+!(list.size()==0)+ are the most obvious solutions.
\end{Answer}
\begin{Answer}{4.32}
Let $p=$``$x>0$'' and $q=$``$x<y$''.
Then the conditional above can be expressed as $(p\wedge q)\vee (p\wedge \neg q)$.
Using a few of the logical equivalences from Table~\ref{table:logicLaws},
it is not too difficult to see that
 $(p\wedge q) \vee (p\wedge \neg q) = p\wedge (q \vee \neg q) = p\wedge T = p.$
Therefore the code simplifies to:
\begin{code}
    if (x>0) {
        x=y;
    }
\end{code}
\end{Answer}
\begin{Answer}{4.33}
This is not equivalent to the original code.  Consider the case when $x=-1$ and $y=1$, for instance.
\end{Answer}
\begin{Answer}{4.34}
After my head is done spinning trying to understand what they are
saying, I suspect something isn't right here. In fact,
since both conditions need to be true when if statements are nested, it is the same thing as a conjunction.
In other words, the two ifs are equivalent to \verb+if( x>0 && (x<y || x>0) )+.
By absorption, this is equivalent to \verb+if(x>0)+.
So the simplified code is:
\begin{code}
    if(x>0) {
        x=y;
    }
\end{code}
You can also think about it this way.  The assignment \verb+x=y+ cannot happen unless \verb+x>0+ due to the outer
if.\footnote{If you think about it, this is why the solution to this in the previous example failed.}
But the inner if has a disjunction, one part of which is \verb+x>0+, which we already know is true.
In other words, it doesn't matter whether or not \verb+x<y+.  This argument also leads to the solution we just gave.
\end{Answer}
\begin{Answer}{4.36}
$p$ is evaluated. If it is false, then the whole expression is evaluated
as false without evaluating $q$ since no matter what the truth
value of $q$ is, $p\wedge q=F$. On the other hand, if
$p$ is true, then the truth value of $p\wedge q$ is the same as
the truth value of $q$, so it evaluates $q$ and uses that as
the value of the whole expression.
\end{Answer}
\begin{Answer}{4.37}
\begin{solevalenum}
\item
This solution is incorrect.
There are a few problems.  The obvious one is that the first statement actually prevents the program
from crashing so it is certainly not unnecessary!  Also, the second and third statements may be equivalent,
but how are they connected?  For instance, given the expression $\neg (A\wedge B)\vee \neg A$,
I cannot simply remove the `redundant' $\neg A$ to obtain an ``equivalent'' expression of $\neg (A\wedge B)$
(if necessary, plug in different truth values for $A$ and $B$ to convince yourself that these are not the same).
\item
This is not correct.  The second part of the expression seems to have disappeared.
But how can we know it isn't equivalent?  We just need to find a scenario where the two versions do different
things. Notice that the `simplified' expression is true when the list is not empty regardless of the value of
element $0$.  But what if the list is not empty and element $0$ is 50?  The original expression is false and the
`simplified' expression is true.  Clearly not the same.
\item
This solution is not only correct, but it is very well argued.
\end{solevalenum}
\end{Answer}
\begin{Answer}{4.38}
Technically speaking, the final solution is not equivalent.
However, it turns out that it is {\em better} than the
original.  This is because the original code would actually crash if the list is empty.
Go back and look at the code and verify that this is the case.
Then verify that the
final simplified version will {\em not} crash.
\end{Answer}
\begin{Answer}{4.43}
\verb+11110000+; \verb+11110000+; \verb+00001111+; $15$.
\end{Answer}
\begin{Answer}{4.44}
\verb+00111001+.
\end{Answer}
\begin{Answer}{4.47}
 \verb+11000000+; \verb+11111100+; \verb+00111100+.
 
\end{Answer}
\begin{Answer}{4.50}
It will return 1 for negative numbers which does not really make sense.
%It will loop forever if $n<0$.
This should be fixed in two ways.
Since $n!$ is undefined when $n<0$, it can't return the correct answer for negative values.
So the first change is to have it return $-1$ if $n<0$.
Some other value can be used, but $-1$ works well because $n!$ can't be negative, so if it returns $-1$, you know something is up.
Second, this behavior should be clearly documented so that it clearly states that it returns $n!$ for $n\geq 0$, and $-1$ for $n<0$.
\end{Answer}
\begin{Answer}{4.51}
\begin{solevalenum}
\item
This algorithm always returns $0$.  If you don't see it right away, carefully work through the algorithm
with a few values of $n$.
\item
This is correct.  It doesn't multiply by $1$, but that doesn't change the answer.
\item
This is also correct.  It is just multiplying the values in the reverse order of the other examples.
\item
This is incorrect.  It actually computes $(n-1)!$.
To fix this, does $i$ have to start at $0$, or go to $n$ (instead of $n-1$)?
We'll leave it to you to work out which is correct.
\end{solevalenum}
\end{Answer}
\begin{Answer}{4.52}
This isn't really that much different than the algorithms to compute $n!$.
Here is one algorithm that does the job:
\begin{code}
    double power(double x, int n) {
        power = 1;
        for(int i=0;i<n;i++) {
            power = power*x;
        }
        return power;
    }
\end{code}
The loop could also have been \verb@for(int i=1;i<=n;i++)@,
or any loop that executes $n$ times since the loop index, $i$, is not used as part
of the calculation.
\end{Answer}
\begin{Answer}{4.55}
Here is one possible solution.  Note that the parentheses are necessary.
\begin{code}
    boolean startsOrEndsWithZero(int[] a, int n) {
        if( n>0 && (a[0]==0 || a[n-1]==0) ) {
            return true;
        else {
            return false;
        }
    }
\end{code}
\end{Answer}
\begin{Answer}{4.56}
The solution uses the expression  \verb+n>0 && (a[0]==0 || a[n-1]==0)+.
If $n=0$, the expression is false because of the \verb+&&+, so the algorithm returns false as it should
since an array with no elements certainly does not begin or end with a $0$.
If $n=1$, first note that $n-1=0$, so $a[0]$ and $a[n-1]$ refer to the same element.
Although this is redundant, it isn't a problem.
If $a[0]=0$, the expression evaluates to $T\wedge (T\vee T)=T$, and the algorithm returns
true as expected.
If $a[0]\neq 0$, the expression evaluates to $T\wedge (F\vee F)=F$, and the algorithm returns
false as expected.
\end{Answer}
\begin{Answer}{4.59}
As we already discussed, many languages truncate when performing integer division.
When the numbers are positive (as they are here), that is the same thing as taking the floor.
Even if a language does not do this, Theorem~\ref{floorThm} implies that it would still work.
\end{Answer}
\begin{Answer}{4.60}
It is easiest to see that this is correct by comparing with the previous solution.
The only difference is that the condition went from $i<=(n-2)/2$ to $i<n/2$.
Notice that $(n-2)/2+1=n/2-2/2+1=n/2$.
But since we also replaced $<=$ with $<$, it still stops at the same point.
\end{Answer}
\begin{Answer}{4.62}
It is correct.  To convince yourself (but this is {\em not} a proof),
plug the numbers $1$ through $5$ (or some other set of both even and odd values) into both sides
to see that you get the same number.
\end{Answer}
\begin{Answer}{4.65}
They are equivalent by the commutative property.
But they are not practically equivalent because
whenever \verb+x>=a.length+, the first statement works
properly because it does not try to access \verb+a[x]+
which avoids an \verb+IndexOutOfBoundsException+,
but the second one will try to access \verb+a[x]+ first,
causing an \verb+IndexOutOfBoundsException+ {\em before}
it checks if the index is valid.
\end{Answer}
\begin{Answer}{4.66}
 This is simple---we just need to add a check
that the index is not negative:\\
\verb+if(x>=0 && x<a.length && a[x]!=0)+.
\end{Answer}
\begin{Answer}{4.69}
Using De Morgan's Law, we get \verb+!(P(i) || Q(i))+.
\end{Answer}
\begin{Answer}{4.70}
First notice that if $P(i)$ is true for every value of $i$, \verb+result+ will be true at the
end of the first loop, so  \verb+isTrueForAll3+ will return true without even considering $Q$.
However, if $P(i)$ is false for any value of $i$, then it will go onto the second loop.
The second loop will return false if $Q(i)$ is false for any value of $i$.  But if $Q(i)$ is true
for all values of $i$, the method returns true.
So, how do we put this all together into a simple answer?
Notice that the only time it returns true is if either $P(i)$ is always true or if $Q(i)$ is always true.
In other words,  \verb+isTrueForAll3+ is determining the truth value of $(\forall i P(i)) \vee (\forall i Q(i))$.
By the way, notice that I used the variable \verb+i+ for both quantifiers. This is sort of like using
the same variables for two separate loops.
\end{Answer}
\begin{Answer}{4.72}
\begin{code}
boolean has2MultipleOr3Multiple(int[] a, int n) {
   for(int i=0;i<n;i++) {
       if(a[i]%2==0 || a[i]%3==0) {
           return true;
       }
   }
    return false;
}
\end{code}
\end{Answer}
\begin{Answer}{4.76}
Here is one answer:\index{sequential search}
\begin{code}
int sequentialSearch(int []a, int n, int val) {
    int i=0;
    while(i<n) {
       if( a[i]==val ) {
           return i;
       }
       i++;
   }
   return -1;
\end{code}
Notice that this algorithm is almost identical to the code from
the previous example---I just made 3 minor changes to the code!
This is because this is using a pattern that is common when
dealing with arrays.
\end{Answer}
\begin{Answer}{4.78}
 In the first iteration in the loop,
$n = 5, i = 1$, $n*i > 4$ and thus $x = 10.$
Next, $n = 3$, $i = 2$, and we go through the loop again.
Since $n*i > 4$, $x = 10 + 2*3 = 16$.
Finally, $n = 1, i = 3$, and the loop stops. Hence $x = 16$ is
returned.
\end{Answer}
\begin{Answer}{4.81}
$\lfloor \sqrt{101}\rfloor=10$.  The only primes less than $10$ are $2$, $3$, $5$, and $7$.
Since $101\bmod 2=1$, $101\bmod 3 = 2$, $101\bmod 5 = 1$, and $101\bmod 7 = 3$, none of
which are $0$, Theorem~\ref{thm:erathostenes_primality} tells us that $101$ is prime.
\end{Answer}
\begin{Answer}{4.82}
$323=17*19$, so it is not prime.  I determined this by seeing if any of the primes
no greater than $\lfloor\sqrt{323}\rfloor=17$ were factors.
Although $2$, $3$, $5$, $7$, $11$, and $13$, are not, $17$ is.
\end{Answer}
\begin{Answer}{4.84}
No need to test the even integers. So
we can modify the code to skip checking even numbers,
so long as we first make sure that it gets the correct
answer for 2. If you didn't think about this,
go back and try to write the algorithm before
continuing to read the solution.
Here is the modified code:
\begin{code}
boolean isPrime(int n) {
    if(n<=1) {      // Anything less than 2 is not prime
        return false;
    } else if (n==2) {
       return true;    // Need to deal with this as a special case
    } else if( n%2==0 ) {  // Discard even numbers.
        return false;
    } else {
        // Determine if it has any odd factors.
        int i = 3;
        while( i <= sqrt(n) ) {
            if( n%i==0 ) {
                return false;
            }
            i = i + 2;
        }
        return true;  // It had no factors.
    }
}
\end{code}
\end{Answer}
\begin{Answer}{4.85}
The following algorithm does the job.
\begin{code}
	int reverseDigits(int n) {
		x=0;
		while(n!=0) {
		    x = x*10+n%10;
		    n=n/10;
		}
		return x;
	}
\end{code}
\end{Answer}
\begin{Answer}{4.88}
Below is a table of the values the algorithm computes.
The final answer of 8388608 is correct!
It took 9 steps as opposed to the 22 required by the
algorithm from Example \ref{exa:powers1}, so it was
definitely faster. Notice that
$2 \log_2(23) \approx 2*4.52356=9.04712>9$,
it works as advertised.

\begin{tabular}{|l|l|l|}
\hline
k&pow&ans\\
\hline
23&2.0&1.0\\
22&2.0&2.0\\
11&4.0&2.0\\
10&4.0&8.0\\
5&16.0&8.0\\
4&16.0&128.0\\
2&256.0&128.0\\
1&65536.0&128.0\\
0&65536.0&8388608.0\\
\hline
\end{tabular}
\end{Answer}
